\documentclass[11pt]{article}

\usepackage[a4paper,margin=2.5cm]{geometry}
\usepackage{booktabs}
\usepackage{caption}
\usepackage{helvet}
\renewcommand{\familydefault}{\sfdefault}
\captionsetup{labelfont=bf,font=small}
\setlength{\parindent}{0pt}
\setlength{\parskip}{6pt}

\renewcommand{\thetable}{S\arabic{table}}
\renewcommand{\thefigure}{S\arabic{figure}}

\title{\vspace{-1.5cm}\Large Supplementary Materials\\[4pt]
\large Interpretable Machine Learning for One-Part Fly-Ash/Slag Geopolymer
Strength Prediction: Toward Multifunctional Binder Design}
\author{Vinoth Nageshwaran, Sudhir Amritphale, and Soundararajan Ezekiel}
\date{}

\begin{document}
\maketitle

This document provides the supplementary tables referenced in the main text.

\vspace{0.5cm}

%============================ TABLE S1 ============================
\begin{table}[h]
\centering
\caption{Bibliometric screening and tagging summary (PRISMA-style). Records were
retrieved from the Crossref REST API via eight bibliographic queries, pooled, and
de-duplicated by DOI. Relevance and sub-stream tagging used case-insensitive
regular expressions applied to record \emph{titles} (the corpus retains
title-level metadata; abstracts were not stored).}
\label{tab:s1}
\begin{tabular}{lr}
\toprule
\textbf{Screening / tagging stage} & \textbf{Count} \\
\midrule
Records retrieved (8 Crossref queries, pooled) & $\approx$6{,}400 (with duplicates) \\
Unique records after DOI de-duplication & 4{,}312 \\
On-topic (geopolymer / alkali-activated regex), 2014--mid-2026 & 1{,}850 \\
\quad--- tagged ML stream & 346 \\
\quad--- tagged one-part stream & 308 \\
\quad--- one-part AND ML & 17 \\
Records with title text $>$120 characters (concept network) & 1{,}092 \\
Validation sample (title-based adjudication) & 100 \\
\bottomrule
\end{tabular}
\end{table}

Regex-tag validation on the 100-record adjudication sample gave precision $=1.00$
for both the ML and one-part streams (no false positives), with recall $=0.76$
(ML) and $0.86$ (one-part); the residual errors are recall misses on records whose
titles omit the trigger terms. Because tagging is title-based, this validation is
conservative on recall.

\vspace{0.5cm}

%============================ TABLE S2 ============================
\begin{table}[h]
\centering
\caption{Leave-one-source-out (LOSO) fold-by-fold results for the XGBoost model on
the 80-mix one-part fly-ash/GGBS geopolymer dataset. Each row holds out all mixes
from one source study and predicts them from a model trained on the other eleven.
Fold $R^2$ on a handful of held-out points is unstable and ranges from strongly
negative to near-unity; per-fold RMSE is a steadier per-study read. The pooled LOSO
$R^2$ over all 80 predictions is 0.61 (five-fold CV mean $R^2 = 0.65$).}
\label{tab:s2}
\begin{tabular}{ccrr}
\toprule
\textbf{Source study (ref.)} & \textbf{Mixes held out ($n$)} & \textbf{LOSO $R^2$} & \textbf{LOSO RMSE (MPa)} \\
\midrule
17 & 6  & 0.097  & 14.35 \\
18 & 6  & -0.154 & 12.45 \\
19 & 4  & -2.196 & 21.80 \\
20 & 6  & -2.720 & 9.33 \\
21 & 4  & 0.673  & 6.94 \\
22 & 2  & -2.584 & 6.06 \\
23 & 3  & 0.996  & 1.01 \\
24 & 6  & 0.474  & 13.07 \\
25 & 14 & -0.349 & 9.30 \\
26 & 11 & 0.852  & 6.13 \\
27 & 7  & 0.060  & 25.26 \\
28 & 11 & -0.218 & 25.61 \\
\midrule
\textbf{Total / pooled} & \textbf{80} & \bfseries 0.612 & \bfseries 15.49 \\
\bottomrule
\end{tabular}
\end{table}

Source-study reference numbers correspond to the citation keys in the main-text
reference list. The wide spread of per-fold $R^2$ (from $-2.72$ to $0.996$) reflects
the small size of individual held-out groups ($n = 2$ to $14$); the pooled LOSO
$R^2$ and the five-fold cross-validated mean are the meaningful aggregate measures
of generalization.

\end{document}
